\section{结论与展望}\label{ch:conclusion}

\subsection{主要结论}

【结论应直接回应研究内容，不要写成摘要的缩写。建议 3--4 条，每条包含“做了什么”和“得到什么结果”。】

（1）本文构建了【数据集/实验流程】，完成了【关键处理步骤】，为【研究问题】提供了可复现的数据基础。

（2）本文提出或采用了【方法名称】，通过【特征/模型/实验设计】实现了【目标任务】，在【测试集/验证集】上取得【指标】。

（3）对比实验和消融实验表明，【关键组件或变量】对性能提升具有主要贡献，说明【机制或应用意义】。

\subsection{创新点}

【本科论文的创新点可以是数据、方法、流程或应用层面的局部创新。避免夸大为“首次”“国际领先”，除非有充分证据。】

（1）数据层面：【填写数据组织、标注、配准或质量控制上的改进】。

（2）方法层面：【填写模型、特征、实验设计或评价方式上的改进】。

（3）应用层面：【填写该方法对实际问题的支撑价值】。

\subsection{不足与展望}

【本节应具体，不要只写“样本量较少、模型有待提高”。建议说明不足如何影响结论，以及下一步怎么解决。】

本文仍存在以下不足：首先，【数据覆盖范围】有限，可能影响模型在其他场景下的泛化能力；其次，【模型或假设】仍较简化，尚未充分考虑【因素】；最后，【验证方式】仍需在独立数据上进一步检验。

后续研究可从三个方面推进：第一，扩充【时间/空间/对象】维度的数据集；第二，引入【更强模型或机理约束】提升解释性与稳定性；第三，面向实际应用构建自动化处理流程，提高方法可复现性和部署效率。
